\chapter{Marco teórico}

\section{Fundamentos del desempeño ambiental y energético}

La integración del desempeño energético en el desarrollo del sector de la construcción resulta crucial para lograr una arquitectura sostenible. Esto exige soluciones integrales que se implementen desde las fases iniciales de diseño y se respalden durante el desarrollo del proyecto y su etapa operativa. A pesar de los avances tecnológicos alcanzados, los edificios siguen siendo responsables de una parte significativa de las emisiones globales de gases de efecto invernadero, por lo que la optimización del diseño arquitectónico constituye una línea de acción relevante para reducir impactos ambientales.

En este contexto, las herramientas de diseño en etapas tempranas permiten mejorar el desempeño de los proyectos de edificación. Ajustes como la optimización de vanos, la configuración de losas, el control de envolventes o la selección de materiales pueden incidir en la reducción del carbono incorporado y en la mejora de la eficiencia energética. En consecuencia, el diseño temprano y los materiales innovadores no solo reducen impactos ambientales, sino que redefinen el papel de la construcción como un vehículo hacia la sostenibilidad.

\section{Impacto ambiental de los edificios}

El impacto ambiental de los edificios está relacionado con el diseño, los materiales, los sistemas técnicos y los patrones de ocupación previstos. La ocupación y el uso del edificio influyen en la demanda energética, la ventilación, las cargas internas y las condiciones de confort. Por ello, el desempeño energético requiere un enfoque integral que articule decisiones de diseño, operación y comportamiento de los usuarios.

En edificaciones administrativas, la reducción del consumo energético debe equilibrarse con la necesidad de mantener condiciones adecuadas de confort térmico, lumínico y ambiental. La gestión del edificio, los horarios de uso y la configuración de los sistemas activos y pasivos son variables que condicionan su comportamiento durante la etapa operativa.

\section{Sostenibilidad en la arquitectura}

La sostenibilidad en la arquitectura articula objetivos ambientales, sociales y económicos a lo largo del ciclo de vida de la edificación. No solo exige la reducción de emisiones y consumos operativos, sino que también considera la huella de carbono de los materiales, el confort de los ocupantes y la resiliencia frente al cambio climático.

Orientar la sostenibilidad hacia su aplicación práctica permite incorporar simulaciones energéticas en etapas tempranas de diseño, reducir la demanda energética, optimizar la envolvente y seleccionar materiales con menor impacto ambiental. Estas herramientas justifican técnicamente las decisiones proyectuales y facilitan la adopción de políticas de construcción sostenible en el sector público.

\section{Eficiencia energética en edificaciones}

El análisis de la eficiencia energética en edificaciones permite evaluar el desempeño funcional de los materiales y la efectividad de las estrategias de control ambiental orientadas a reducir el impacto energético durante el ciclo de vida del edificio. Este ciclo comprende las fases de diseño, implementación, operación y eventual desmantelamiento.

La eficiencia energética debe entenderse como una respuesta al cambio climático, en la que las edificaciones actúan como sistemas integrados que combinan materiales, sistemas constructivos, análisis de ciclo de vida y estrategias de operación. Su objetivo principal es minimizar el impacto ambiental sin comprometer la habitabilidad ni el confort interior.

\section{Estrategias pasivas de diseño}

Las estrategias pasivas constituyen un eje fundamental en el diseño sostenible de las edificaciones. El aprovechamiento de factores climáticos, como radiación solar, vientos dominantes, sombra, inercia térmica y ventilación natural, permite reducir la dependencia de sistemas mecánicos y optimizar el consumo energético durante la operación del edificio.

\subsection{Orientación y control solar}

La orientación y el control solar forman parte de las estrategias principales del diseño pasivo. La correcta interacción entre radiación solar, ganancia térmica e iluminación natural permite reducir cargas energéticas y mejorar las condiciones de confort. En climas cálidos, estas decisiones deben adoptarse desde etapas tempranas del proyecto para evitar sobrecalentamiento y reducir la dependencia de sistemas activos.

\subsection{Ventilación natural y confort térmico}

En edificios administrativos, uno de los principales retos consiste en alcanzar condiciones adecuadas de confort térmico sin recurrir de manera intensiva a sistemas activos de climatización. La ventilación natural, la ventilación cruzada, la ventilación nocturna y el control solar pueden reducir la carga térmica interna y mejorar el confort adaptativo en climas cálidos.

Estos criterios resultan pertinentes para contextos insulares como Galápagos, donde las condiciones ambientales demandan soluciones bioclimáticas de bajo impacto. La integración de estrategias pasivas puede reducir la carga térmica, mejorar la calidad ambiental interior y disminuir el consumo energético asociado a la climatización.

\subsection{Iluminación natural y eficiencia lumínica}

La iluminación natural es una estrategia pasiva clave para la eficiencia energética y el control lumínico de las edificaciones. Su correcta aplicación permite reducir la dependencia de iluminación artificial y mejorar la calidad ambiental de los espacios interiores. En contextos tropicales, debe responder simultáneamente a la necesidad de aprovechar la luz natural y mitigar el deslumbramiento.

\subsection{Aislamiento y envolvente térmica}

El aislamiento térmico y la configuración de la envolvente son factores determinantes en el comportamiento energético de una edificación. En climas tropicales, la radiación solar y las altas temperaturas condicionan el desempeño térmico, por lo que la selección del material, el espesor, la masa térmica y la capacidad de almacenamiento de calor deben ser evaluados de forma integrada.

\section{Materiales sostenibles en la construcción}

Los materiales sostenibles constituyen una parte fundamental del desarrollo de nuevas edificaciones, ya que responden a necesidades técnicas, ambientales y territoriales. Su implementación busca minimizar el consumo energético, la generación de residuos y las emisiones asociadas al ciclo de vida de la construcción.

Los materiales sostenibles pueden agruparse en reciclados, renovables y locales. Los materiales reciclados se vinculan con la economía circular y la reducción de energía incorporada; los materiales renovables proceden de recursos que pueden regenerarse en periodos relativamente cortos; y los materiales locales reducen emisiones asociadas al transporte y fortalecen la relación entre arquitectura, territorio y cultura constructiva.

\section{Análisis de ciclo de vida}

El análisis de ciclo de vida constituye un componente esencial en el desarrollo de edificaciones sostenibles, ya que permite evaluar impactos desde la extracción de materiales hasta la operación, mantenimiento y fin de vida del edificio. Este enfoque facilita la comparación entre materiales convencionales y sostenibles, y orienta decisiones de diseño con base en indicadores ambientales cuantificables.

En el contexto de Galápagos, el análisis de ciclo de vida adquiere relevancia por la condición insular, la logística de transporte y la sensibilidad ambiental del territorio. La selección de materiales debe considerar no solo su desempeño térmico, sino también su procedencia, durabilidad, mantenimiento y posibilidad de reutilización.

\section{Normativas y estándares}

La aplicación de estándares internacionales como LEED, EDGE, WELL, CEELA y Minergie, junto con la normativa ecuatoriana y local, permite establecer un marco de referencia para evaluar el desempeño ambiental de las edificaciones. Estos sistemas integran criterios de eficiencia energética, confort, agua, materiales, calidad ambiental interior y reducción de impactos.

Para el caso de San Cristóbal, el proyecto se fundamenta en regulaciones locales orientadas al diseño pasivo, la ventilación natural, el aislamiento térmico y el uso de materiales locales y sostenibles. Estas regulaciones buscan minimizar la demanda energética y el impacto ambiental, priorizando estrategias que respondan a las condiciones climáticas y ecológicas de las islas.

\section{Modelación y simulación energética}

La simulación energética es una herramienta esencial para sustentar decisiones de diseño orientadas a mejorar el comportamiento térmico, el consumo energético y las condiciones de confort. La interoperabilidad entre modelos BIM y motores de cálculo energético, como EnergyPlus, permite integrar el modelado geométrico con la evaluación cuantitativa del desempeño.

En esta investigación, la combinación BIM-BES se plantea como base metodológica para evaluar el escenario actual y comparar alternativas de optimización. Herramientas como Revit y DesignBuilder permiten modelar el edificio, definir propiedades de materiales, simular escenarios y cuantificar indicadores como demanda energética, consumo, confort térmico, ventilación e iluminación natural.
